\clearpage
\pagestyle{fancy}
\markboth{摘要}{摘要}
\phantomsection
\addcontentsline{toc}{section}{摘\quad 要}
\pagenumbering{Roman}
\setcounter{page}{1}

% 标题：摘要
\begin{center}
    \setstretch{1.5}
    \zihao{3}\heiti
    \vspace*{12pt}
    摘\quad\quad 要
    \vspace*{6pt}
\end{center}

% 摘要正文
\begingroup
    \zihao{-4}\songti
    \setstretch{1.25}
    \setlength{\parindent}{2em}
    \setlength{\parskip}{0pt}
    
    本文充分运用煤地球化学、煤岩学、煤田地质学、岩石学和矿物学等理论知识，研究了煤中伴生元素的地质地球化学习性和成因机理，归纳了煤中伴生元素的 7 种成因模式。系统总结和分析了中国华北和西南地区煤中伴生元素的分布特征及其差异的地质成因。详细讨论了煤中稀土元素的赋存分配模式和机理。指出了煤中铂族元素的背景值及 5 种来源，提出煤中铂族元素以 Pt-Pd 分配模式为特征。应用 TOF-SIMS 研究了高硫煤中黄铁矿化的杆状菌落的伴生元素组成，指出菌落在生物成因黄铁矿的形成过程中，在活化和富集 Cu、Zn 等金属元素性研究了同沉积火山灰对煤中伴生元素富集和存在形态的影响，发现了煤中以非硫化物状态存在的高含量 Fe、Cu 等元素，并给予了由火山灰、有机质和陆源碎屑组成特殊结构的岩石命名和分类方案。
    \par
    \vspace{\baselineskip}
    \centerline{\color{red}...... 1500 字左右 ......}
    \par
    \vspace{\baselineskip}
    在风氧化作用对煤中伴生元素富集影响方面，指出 U、Mo、Cu 等在还原障和氧化障中具有高度的迁移反差，在氧化障中具有较高的活性，而在还原障中可形成次生堆积。
    \par
\endgroup

\par
\vspace{6pt}
\begingroup
    \setstretch{1.25}
    \zihao{-4}
    \noindent\hspace{2em}{\heiti 关键词：}{\songti 煤；伴生元素；地质地球化学习性；富集模式}
\endgroup
\clearpage
